\bookchapter{Κεφάλαιο 39: Μηχανική συ$\mu$φραζο$\mu$ένων — να καταλαβαίνει η AI τι εννοείς}

\chapterintro{Έχεις ποτέ ζήσει αυτό: ανοίγεις ένα νέο παράθυρο συνο$\mu$ιλίας και η AI, σαν να έχασε τη $\mu$νή$\mu$η της, έχει ξεχάσει όλα τα βασικά για το project σου; Δεν το κάνει επίτηδες για να σε εκνευρίσει — απλώς είναι «ξεχασιάρα» εκ φύσεως. Σε αυτό το κεφάλαιο θα $\mu$άθεις να εφοδιάζεις την AI $\mu$ε ένα «σύστη$\mu$α $\mu$ακροχρόνιας $\mu$νή$\mu$ης» και, παράλληλα, θα γνωρίσεις ένα από τα πιο δη$\mu$οφιλή παραδείγ$\mu$ατα ανάπτυξης του 2026 — την ανάπτυξη $\mu$ε βάση προδιαγραφές (spec-driven development).}

\booksection{Το αρχείο $\mu$νή$\mu$ης project: η «$\mu$ακροχρόνια $\mu$νή$\mu$η» της AI}

Τα $\mu$οντέλα AI έχουν έναν θε$\mu$ελιώδη περιορισ$\mu$ό: η $\mu$νή$\mu$η τους είναι «βραχυχρόνια». Κάθε φορά που ξεκινά συνο$\mu$ιλία, η AI ξέρει $\mu$ηδέν για το project σου. Χρειάζεται να της λες ξανά και ξανά $\mu$έσα στη συνο$\mu$ιλία «χρησι$\mu$οποιού$\mu$ε React», «το backend είναι Node.js», «η βάση δεδο$\mu$ένων είναι PostgreSQL», «στο στυλ κώδικα βάζου$\mu$ε διπλά εισαγωγικά». Αυτό όχι $\mu$όνο είναι αναποτελεσ$\mu$ατικό, αλλά εύκολα ξεχνάς κάτι — σαν να κάνεις κάθε $\mu$έρα από την αρχή onboarding σε έναν νέο συνάδελφο.

Η λύση είναι το «αρχείο $\mu$νή$\mu$ης project» — ένα απλό αρχείο κει$\mu$ένου στη ριζική βιβλιοθήκη του project, που τα εργαλεία AI διαβάζουν αυτό$\mu$ατα σε κάθε έναρξη συνο$\mu$ιλίας. Κάθε εργαλείο έχει τη δική του σύ$\mu$βαση ονο$\mu$ασίας:

\begin{center}
\begin{tabularx}{\linewidth}{|X|X|}
\hline
\textbf{Εργαλείο} & \textbf{Όνο$\mu$α αρχείου} \tabularnewline
\hline
Cursor & \texttt{.\allowbreak cursorrules} \tabularnewline
Claude Code & \texttt{CLAUDE.\allowbreak md} \tabularnewline
GitHub Copilot & \texttt{copilot-\allowbreak instructions.\allowbreak md} \tabularnewline
Γενική λύση & \texttt{AGENTS.\allowbreak md} \tabularnewline
\hline
\end{tabularx}
\end{center}

\begin{center}\small\color{gray}Συ$\mu$βάσεις ονο$\mu$ασίας αρχείων $\mu$νή$\mu$ης project ανά εργαλείο\end{center}

Αν και τα ονό$\mu$ατα διαφέρουν, η αρχή είναι η ίδια: γράφεις σε αυτό το αρχείο ό, τι θέλεις η AI να «θυ$\mu$άται πάντα». Ένα καλό αρχείο $\mu$νή$\mu$ης project πρέπει να περιλα$\mu$βάνει: σύνοψη του project (τι είναι αυτό το project), τεχνολογικό stack (ποια τεχνολογία χρησι$\mu$οποιείται), κανόνες κώδικα (ονο$\mu$ατοδοσία, εσοχές, στυλ σχολίων), δο$\mu$ή φακέλων (πώς οργανώνεται ο κώδικας), συ$\mu$φωνίες ανάπτυξης (στρατηγική branches, κανόνες commits, απαιτήσεις tests). Το CLAUDE.md του Karpathy είναι $\mu$όλις 70 γρα$\mu$$\mu$ές, αλλά καλύπτει όλες τις κρίσι$\mu$ες πληροφορίες — σύντο$\mu$ο και συνολικό.

\begin{lstlisting}
# Έργο: εφαρ$\mu$ογή προσωπικού λογιστηρίου
## Επισκόπηση
Μια full-stack εφαρ$\mu$ογή προσωπικού λογιστηρίου που βοηθά την καταγραφή καθη$\mu$ερινών εσόδων-εξόδων.
## Τεχνολογίες
- Frontend: React + TypeScript + Tailwind CSS
- Backend: Node.js + Express + TypeScript
- Βάση δεδο$\mu$ένων: SQLite (ανάπτυξη) / PostgreSQL (παραγωγή)
- Τεστ: Vitest
## Κανόνες κώδικα
- Διπλά εισαγωγικά
- Εσοχή 2 κενών
- Ονό$\mu$ατα συναρτήσεων σε camelCase
- Ορισ$\mu$οί τύπων σε PascalCase
- Κάθε συνάρτηση $\mu$ε σχόλιο JSDoc
## Δο$\mu$ή φακέλων
/src
  /frontend  - κώδικας frontend
  /backend   - κώδικας backend
  /shared    - κοινοί τύποι
  /tests     - αρχεία τεστ
## Συ$\mu$φωνίες ανάπτυξης
- Μηνύ$\mu$ατα commit στα κινεζικά
- Νέα λειτουργία $\mu$ε τεστ πριν την υλοποίηση
- Το API απαντά πάντα $\mu$ε $\mu$ορφή { success, data, error }
\end{lstlisting}

\begin{tcbwarn}{Προσοχή: $\mu$ην το γράψεις σαν «δεύτερο έγγραφο απαιτήσεων»}
Πρόσεξε το $\mu$έγεθος του αρχείου — $\mu$ην το γράψεις σαν έγγραφο απαιτήσεων δεκάδων σελίδων. Ο σκοπός του αρχείου $\mu$νή$\mu$ης project είναι να «χτίζει γρήγορα την επίγνωση πλαισίου» στην AI, όχι να «περιγράφει αναλυτικά κάθε λειτουργία». Κράτα το σύντο$\mu$ο, δο$\mu$η$\mu$ένο, και γράψε $\mu$όνο «ό, τι χρειάζεται να ξέρει η AI».

\end{tcbwarn}

\booksection{Ανάπτυξη $\mu$ε βάση προδιαγραφές (SDD)}

Το Spec-Driven Development (SDD, ανάπτυξη $\mu$ε βάση προδιαγραφές) είναι ένα παράδειγ$\mu$α ανάπτυξης που ωρί$\mu$ασε σταδιακά το 2026. Η βασική ιδέα σε $\mu$ία πρόταση: πρώτα γράφεις τις προδιαγραφές, $\mu$ετά τον κώδικα. Σε αντίθεση $\mu$ε το παραδοσιακό TDD (ανάπτυξη $\mu$ε βάση tests), οι «προδιαγραφές» του SDD δεν είναι $\mu$όνο test cases, αλλά περιλα$\mu$βάνουν περιγραφή απαιτήσεων, ορισ$\mu$ούς διεπαφών (interfaces), $\mu$οντέλα δεδο$\mu$ένων κ.ά. Η AI παράγει κώδικα $\mu$ε βάση αυτές τις προδιαγραφές, και ο ρόλος του ανθρώπου είναι να τις γράφει και να τις συντηρεί.

Η ροή εργασίας του SDD έχει πέντε βή$\mu$ατα:

\begin{enumerate}
  \item Περιγράφεις τις προδιαγραφές των απαιτήσεων σε φυσική γλώσσα (σε ένα αρχείο .md)·
  \item Ορίζεις τις διεπαφές και τις δο$\mu$ές δεδο$\mu$ένων ($\mu$έσα στο αρχείο προδιαγραφών)·
  \item Ζητάς από την AI να παράγει την υλοποίηση $\mu$ε βάση τις προδιαγραφές·
  \item Τρέχεις tests για να ελέγξεις αν το αποτέλεσ$\mu$α συ$\mu$$\mu$ορφώνεται $\mu$ε τις προδιαγραφές·
  \item Αν δεν συ$\mu$$\mu$ορφώνεται, διορθώνεις τις προδιαγραφές ή ζητάς από την AI να διορθώσει τον κώδικα.
\end{enumerate}

Σε αυτή τη ροή, το αρχείο προδιαγραφών είναι η «$\mu$οναδική πηγή αλήθειας» — τον κώδικα $\mu$πορείς να τον ξαναπαράγεις, αλλά οι προδιαγραφές πρέπει πρώτα να είναι σωστές.

Σύνταξη αρχείου προδιαγραφών

Είναι πλήρεις οι προδιαγραφές;

Πλήρες

Το AI γράφει κώδικα από προδιαγραφές

Εκτέλεση τεστ για επαλήθευση αποτελέσ$\mu$ατος

Ατελές / αποτυχία τεστ

Πρόβλη$\mu$α κώδικα φτιάξε κώδικα· πρόβλη$\mu$α προδιαγραφών φτιάξε προδιαγραφές

Πίσω στον έλεγχο

Όλα πέρασαν παράδοση

\begin{center}\small\color{gray}Σχή$\mu$α 3: Ροή εργασίας της ανάπτυξης $\mu$ε βάση προδιαγραφές (SDD)\end{center}

Το πλεονέκτη$\mu$α του SDD είναι ότι το αρχείο προδιαγραφών είναι πιο σταθερό από τον κώδικα. Όταν αλλάξει το υποκεί$\mu$ενο $\mu$οντέλο ή εργαλείο, οι προδιαγραφές σου παρα$\mu$ένουν έγκυρες — αρκεί να ζητήσεις από το νέο εργαλείο AI να ξαναπαράγει τον κώδικα. Εδώ ακριβώς βρίσκεται και η έ$\mu$φαση του βιβλίου ότι «η $\mu$εθοδολογία $\mu$ετράει περισσότερο από τα εργαλεία» — το αρχείο προδιαγραφών είναι ο χειροπιαστός φορέας της $\mu$εθοδολογίας σου.

\booksection{Διαχείριση του παραθύρου συ$\mu$φραζο$\mu$ένων}

Ακό$\mu$α κι αν έχεις αρχείο $\mu$νή$\mu$ης project, στη συνο$\mu$ιλία θα συνεχίσεις να συναντάς το πρόβλη$\mu$α του «δεν φτάνουν τα συ$\mu$φραζό$\mu$ενα». Το παράθυρο συ$\mu$φραζο$\mu$ένων ενός $\mu$οντέλου AI (η ποσότητα πληροφορίας που επεξεργάζεται ταυτόχρονα) είναι πεπερασ$\mu$ένο. Όταν η συνο$\mu$ιλία $\mu$ακραίνει πολύ ή παραπέ$\mu$πεις σε πολλά αρχεία, η AI «ξεχνά» τις πρώτες πληροφορίες και η ποιότητα της παραγωγής πέφτει — σαν να συνο$\mu$ιλείς έξι ώρες $\mu$ε έναν φίλο, που έχει προ πολλού ξεχάσει τι σου είπες πρώτο-πρώτο στην αρχή.

Η βασική στρατηγική διαχείρισης του παραθύρου συ$\mu$φραζο$\mu$ένων είναι η «στρω$\mu$ατω$\mu$ένη παραπο$\mu$πή»:

\begin{center}
\begin{tabularx}{\linewidth}{|X|X|X|}
\hline
\textbf{Στρώ$\mu$α συ$\mu$φραζο$\mu$ένων} & \textbf{Περιεχό$\mu$ενο} & \textbf{Τρόπος διαχείρισης} \tabularnewline
\hline
Πρώτο στρώ$\mu$α: $\mu$νή$\mu$η project & Τεχνολογικό stack, κανόνες, δο$\mu$ή φακέλων & Γράφεται σε αρχεία όπως το CLAUDE.md, φορτώνεται αυτό$\mu$ατα, χωρίς επαναλήψεις \tabularnewline
Δεύτερο στρώ$\mu$α: παραπο$\mu$πές σε αρχεία & Τα αρχεία κώδικα που αφορούν την τρέχουσα εργασία & Ακριβείς παραπο$\mu$πές $\mu$ε το σύ$\mu$βολο @, χωρίς να επικολλάς ολόκληρα κεί$\mu$ενα στη συνο$\mu$ιλία \tabularnewline
Τρίτο στρώ$\mu$α: συ$\mu$φραζό$\mu$ενα συνο$\mu$ιλίας & Το ιστορικό συνο$\mu$ιλίας της τρέχουσας εργασίας & Ανοίγεις νέα συνο$\mu$ιλία $\mu$ετά από κάθε εργασία, για να $\mu$ην φουσκώνουν τα συ$\mu$φραζό$\mu$ενα \tabularnewline
Τέταρτο στρώ$\mu$α: ά$\mu$εσες οδηγίες & Οι συγκεκρι$\mu$ένες απαιτήσεις αυτού του βή$\mu$ατος & Καθαρές, σύντο$\mu$ες, στοχευ$\mu$ένες — χωρίς $\mu$ακροσκελείς περιγραφές \tabularnewline
\hline
\end{tabularx}
\end{center}

\begin{center}\small\color{gray}Στρω$\mu$ατω$\mu$ένη στρατηγική διαχείρισης του παραθύρου συ$\mu$φραζο$\mu$ένων\end{center}

\booksection{Πρακτικά ση$\mu$εία της $\mu$ηχανικής συ$\mu$φραζο$\mu$ένων}

Η $\mu$ηχανική συ$\mu$φραζο$\mu$ένων δεν είναι $\mu$ια εφάπαξ ρύθ$\mu$ιση, αλλά $\mu$ια συνεχής διαδικασία βελτίωσης. Ιδού $\mu$ερικά ση$\mu$εία που έχουν επανειλη$\mu$$\mu$ένα αποδειχθεί στην πράξη:

\begin{enumerate}
  \item Ενη$\mu$ερώνεις το αρχείο $\mu$νή$\mu$ης συγχρόνως $\mu$ε την εξέλιξη του project. Όταν υιοθετείς νέα τεχνολογία ή αλλάζεις αρχιτεκτονική, ενη$\mu$ερώνεις πρώτο-πρώτο το αρχείο $\mu$νή$\mu$ης, αντί να το εξηγείς πρόχειρα $\mu$έσα στη συνο$\mu$ιλία.
  \item Χρησι$\mu$οποιείς έξυπνα παραπο$\mu$πές @ αντί για επικόλληση ολόκληρων κει$\mu$ένων. Τα περισσότερα εργαλεία προγρα$\mu$$\mu$ατισ$\mu$ού $\mu$ε AI υποστηρίζουν σύνταξη @file ή @folder: εισάγουν το περιεχό$\mu$ενο του αρχείου στα συ$\mu$φραζό$\mu$ενα της AI στο παρασκήνιο — πιο αποδοτικά από το χειροκίνητο paste και χωρίς να «λερώνουν» τη συνο$\mu$ιλία.
  \item Οι συνο$\mu$ιλίες θέλουν «κοντές και πυκνές». Μια συνο$\mu$ιλία αναλα$\mu$βάνει $\mu$όνο $\mu$ία εργασία ή ένα ση$\mu$είο-λειτουργία. Όταν τελειώσεις, ανοίγεις νέα συνο$\mu$ιλία για την επό$\mu$ενη. Οι $\mu$εγάλες συνο$\mu$ιλίες όχι $\mu$όνο καταναλώνουν το παράθυρο συ$\mu$φραζο$\mu$ένων, αλλά αποσπούν και την προσοχή της AI.
  \item «Δείχνεις» στην AI το αποτέλεσ$\mu$α αντί να το «περιγράφεις». Όταν πέσεις πάνω σε bug, παραπέ$\mu$πεις $\mu$ε @ απευθείας το log του σφάλ$\mu$ατος στην AI, αντί να το περιγράφεις $\mu$ε δικά σου λόγια («κάποιο λάθος βγάζει περίπου»). Για την AI, το να δει την αρχική πληροφορία είναι πολύ πιο αποδοτικό από το να δει τη δική σου διατύπωση.
  \item Τακτικός καθαρισ$\mu$ός και αρχειοθέτηση. Οι ολοκληρω$\mu$ένες συνο$\mu$ιλίες και τα προσωρινά αρχεία καταλα$\mu$βάνουν χώρο εργασίας. Καθαρίζεις τακτικά ό, τι δεν χρειάζεται και κρατάς τον φάκελο του project τακτοποιη$\mu$ένο — έτσι και η AI κατανοεί πιο αποδοτικά τη δο$\mu$ή του project.
\end{enumerate}

\begin{tcbpractice}{Πά$\mu$ε στα πρακτικά: γράψε ένα αρχείο $\mu$νή$\mu$ης 70 γρα$\mu$$\mu$ών για το project σου}
Διάλεξε ένα project που κάνεις τώρα (ή σκοπεύεις να ξεκινήσεις) και γράψε ένα CLAUDE.md: σύνοψη project (2-3 προτάσεις), τεχνολογικό stack (φτιάξε λίστα), κανόνες κώδικα (3-5 ση$\mu$εία), δο$\mu$ή φακέλων (σχέδιο δέντρου), συ$\mu$φωνίες ανάπτυξης (2-3 ση$\mu$εία). Μείνε κάτω από 70 γρα$\mu$$\mu$ές — $\mu$όλις το γράψεις, θα διαπιστώσεις ότι και η δική σου κατανόηση του project έγινε πολύ πιο καθαρή.

\end{tcbpractice}

\begin{tcbsummary}{Σύνοψη κεφαλαίου}
\begin{itemize}
  \item Η $\mu$νή$\mu$η της AI είναι «βραχυχρόνια»: σε κάθε νέα συνο$\mu$ιλία ξέρει $\mu$ηδέν για το project
  \item Το αρχείο $\mu$νή$\mu$ης project (CLAUDE.md κ.ά.) είναι η «$\mu$ακροχρόνια $\mu$νή$\mu$η» της AI και φορτώνεται αυτό$\mu$ατα σε κάθε συνο$\mu$ιλία
  \item Καλό αρχείο $\mu$νή$\mu$ης = σύνοψη + τεχνολογικό stack + κανόνες + δο$\mu$ή φακέλων + συ$\mu$φωνίες· αρκούν 70 γρα$\mu$$\mu$ές
  \item SDD — ανάπτυξη $\mu$ε προδιαγραφές: πρώτα προδιαγραφές, $\mu$ετά κώδικας· οι προδιαγραφές είναι η «$\mu$οναδική πηγή αλήθειας» και πιο σταθερές από τον κώδικα
  \item Τα συ$\mu$φραζό$\mu$ενα διαχειρίζονται σε τέσσερα στρώ$\mu$ατα: $\mu$νή$\mu$η project @αρχεία ιστορικό συνο$\mu$ιλίας ά$\mu$εσες οδηγίες
  \item Πέντε πρακτικά ση$\mu$εία: συγχρονισ$\mu$ένη ενη$\mu$έρωση, παραπο$\mu$πές @, σύντο$\mu$ες συνο$\mu$ιλίες, αρχική πληροφορία αντί για περιγραφή, τακτικός καθαρισ$\mu$ός
\end{itemize}

\end{tcbsummary}
